\subsection{theoryv3 Evidence Layer (Compressed, Auto-Generated)}
\label{sec:theoryv3-evidence}

The \texttt{theoryv3} analysis branch is a compact ``evidence layer'' for the discrete kernel and crosslink logic. Current inventory (auto-generated from \codepath{tfpt-suite/2.5/theoryv3/out/}):
\[
\text{modules}=\TFPTTheoryVThreeModulesTotal,\qquad
\text{checks}=\TFPTTheoryVThreeChecksTotal\ (\mathrm{PASS/WARN/FAIL}=\TFPTTheoryVThreeChecksPassTotal/\TFPTTheoryVThreeChecksWarnTotal/\TFPTTheoryVThreeChecksFailTotal),\qquad
\text{plots expected/present}=\TFPTTheoryVThreePlotsExpectedTotal/\TFPTTheoryVThreePlotsPresentTotal.
\]

% Inlined (standalone build): formerly \input{../tfpt-suite/2.5/theoryv3/out/theoryv3_manifest_table.tex}
% -----------------------------------------------------------------------------
% AUTO-GENERATED by tfpt-suite/2.5/theoryv3/theoryv3_manifest.py — do not edit by hand.
\begin{table}[H]
\centering
\small
\begin{tabular}{@{}lllcc@{}}
\toprule
Module & Signal & Result & Checks & Plots OK\\
\midrule
\texttt{alpha\_backreaction\_sensitivity\_audit} & $\alpha$ backreaction sensitivity (k-grid) & PASS & 2/2 & yes\\
\texttt{axion\_dm\_audit} & Axion DM audit & PASS & 2/2 & yes\\
\texttt{baryon\_consistency\_audit} & Baryon consistency / H0 audit & WARN & 3/4 & yes\\
\texttt{constant\_factory\_audit} & Constant factory audit (derived outputs) & WARN & 2/3 & yes\\
\texttt{dark\_energy\_exponential\_audit} & Dark energy exponential candidate & PASS & 2/2 & yes\\
\texttt{dark\_energy\_norm\_half\_origin\_audit} & Dark energy normalization n=1/2 scan & PASS & 2/2 & yes\\
\texttt{defect\_partition\_g5\_audit} & Two-defect partition ($\delta_2=(5/4)\,\delta_{\rm top}^2$) & PASS & 4/4 & yes\\
\texttt{flavor\_pattern\_audit} & Flavor pattern audit (texture invariants) & PASS & 3/3 & yes\\
\texttt{g5\_crosslink\_audit} & Crosslink consistency (g=5 anchors) & PASS & 3/3 & yes\\
\texttt{g5\_origin\_audit} & Holonomy degeneracy g=5 & PASS & 2/2 & yes\\
\texttt{pmns\_tm1\_audit} & PMNS TM1 audit & PASS & 2/2 & yes\\
\texttt{seed\_invariants\_audit} & Seed invariants (c3, $\varphi_0$, $\delta_{\rm top}$, $\beta$) & PASS & 7/7 & yes\\
\texttt{yukawa\_exponent\_index\_audit} & Yukawa exponent-index mapping & WARN & 3/7 & yes\\
\texttt{yukawa\_index\_mapping\_audit} & Yukawa index mapping & WARN & 3/7 & yes\\
\bottomrule
\end{tabular}
\caption{theoryv3 compressed evidence layer (auto-generated from \,\codepath{tfpt-suite/2.5/theoryv3/out/}).}
\end{table}

\begin{tcolorbox}[colback=gray!5, colframe=gray!60,
title={\textbf{Two ``lock-in'' results (negative controls; theoryv3)}}, fonttitle=\bfseries\small]
\textbf{$k=2$ is selected by a sensitivity minimum:} the \texttt{alpha\_backreaction\_sensitivity\_audit} shows the minimum $|\mathrm{ppm}|$ in the $k$ grid occurs at $k=2$, with $|\mathrm{ppm}|\simeq 0.037$.

\textbf{$g=5$ is selected by a defect-partition negative control:} the \texttt{defect\_partition\_g5\_audit} shows $g=5$ minimizes $|z|$ sharply versus $g\in\{4,6\}$ (gap $\sim 45$ in $|z|$).
\end{tcolorbox}

\begin{figure}[H]
\centering
\begin{tikzpicture}
\begin{groupplot}[
    group style={group size=2 by 1, horizontal sep=1.2cm},
    width=0.48\textwidth,
    height=0.42\textwidth,
    grid=both,
    legend style={font=\scriptsize},
]
\nextgroupplot[
    title={$k$-sensitivity (abs ppm)},
    xlabel={$k$},
    ylabel={$|{\rm ppm}|$},
    xmin=-0.1, xmax=3.1,
    ymin=0, ymax=4.2,
    xtick={0,1,1.5,2,2.5,3},
]
\addplot+[mark=*, mark size=2.2pt] coordinates {
    (0.0, 3.665371791638453)
    (1.0, 1.8074246412353618)
    (1.5, 0.883514212976573)
    (2.0, 0.03703710120260405)
    (2.5, 0.9542414840493093)
    (3.0, 1.8681110743677323)
};
\addplot[gray, dashed] coordinates {(2.0,0.0) (2.0,4.2)};

\nextgroupplot[
    title={$g$ negative control (z-score)},
    xlabel={$g$},
    ylabel={$z$},
    xmin=3.7, xmax=6.3,
    ymin=-60, ymax=60,
    xtick={4,5,6},
]
\addplot+[mark=*, mark size=2.2pt] coordinates {
    (4, -46.8455434804275)
    (5, 1.8646873709648437)
    (6, 50.5749164827739)
};
\addplot[gray, dashed] coordinates {(5,-60) (5,60)};
\end{groupplot}
\end{tikzpicture}
\caption{theoryv3 lock-in diagnostics: $k$-sensitivity of $\alpha^{-1}(0)$ (left) and the $g\in\{4,5,6\}$ negative control selecting $g=5$ (right).}
\end{figure}

\begin{tcolorbox}[colback=gray!5, colframe=gray!60,
title={\textbf{Global scorecard (dashboard $\chi^2$; no covariance)}}, fonttitle=\bfseries\small]
The \texttt{global\_consistency\_test} module evaluates a reference-table $\chi^2$ over a small set of headline observables (engineering mode). Current totals:
\[
\chi^2_{\mathrm{core}}=3.083\ (6~\mathrm{dof}),\qquad p\simeq 0.798,
\]
and (excluding $\alpha$ metrology dominance)
\[
\chi^2_{\mathrm{core}\setminus\alpha}=1.354\ (5~\mathrm{dof}),\qquad p\simeq 0.929.
\]
This is a dashboard (independent Gaussians; simplified $r$ bound proxy), not a publication-grade global likelihood. The suite reports per-term contributions and an ``excluding $\alpha$'' total explicitly to avoid misleading optics when a single ultra-precise reference dominates.
\end{tcolorbox}

\begin{figure}[H]
\centering
\begin{tikzpicture}
\begin{groupplot}[
    group style={group size=2 by 1, horizontal sep=1.2cm},
    ybar,
    bar width=10pt,
    width=0.48\textwidth,
    height=0.42\textwidth,
    xmin=0.5, xmax=4.5,
    xtick={1,2,3,4},
    xticklabels={$\beta_{\rm deg}$, $\lambda$, $n_s$, $A_s$},
    x tick label style={rotate=0, font=\scriptsize},
    grid=both,
]
\nextgroupplot[
    title={$\chi^2$ contributions (excluding $\alpha$)},
    ylabel={$\chi^2$},
    ymin=0, ymax=0.8,
]
\addplot coordinates {
    (1,0.5903174784328564)
    (2,0.6542656358088781)
    (3,0.021391549817205795)
    (4,0.08789723472938661)
};

\nextgroupplot[
    title={$z$-scores (excluding $\alpha$)},
    ylabel={$z$},
    ymin=-1.0, ymax=0.15,
]
\addplot coordinates {
    (1,-0.7683212078505034)
    (2,-0.8088668838621581)
    (3,-0.14625850340136054)
    (4,-0.296474678057649)
};
\addplot[black!50, dashed] coordinates {(0.5,0) (4.5,0)};
\end{groupplot}
\end{tikzpicture}
\caption{Engineering-mode dashboard plots from \texttt{global\_consistency\_test}: radar summary (left) and per-observable $\chi^2$ contributions excluding $\alpha$ (right).}
\end{figure}

%=============================================================================
\section{\texorpdfstring{E$_8$ Chain Certificate}{E8 Chain Certificate}}
\label{app:E8cert}
%=============================================================================

This appendix provides the complete verification data for the E$_8$ chain uniqueness 
theorem (Theorem~\ref{thm:chainunique}).

\subsection{Orbit Data Structure}

Each nilpotent orbit $\mathcal{O}$ of $\mathfrak{e}_8$ is characterized by:
\begin{itemize}[leftmargin=1.5em, topsep=2pt, itemsep=1pt]
    \item \textbf{Bala--Carter label:} Standard notation (e.g., $A_1$, $2A_1$, $A_2$, \ldots)
    \item \textbf{Dimension:} $\dim(\mathcal{O})$
    \item \textbf{Centralizer dimension:} $D(\mathcal{O}) = 248 - \dim(\mathcal{O})$
    \item \textbf{Height:} $h(\mathcal{O})$ (sum of weighted Dynkin diagram coefficients)
\end{itemize}

\subsection{The Unique Chain (Excerpt)}

The complete chain from $D = 58$ to $D = 8$:

\begin{center}
\small
\begin{tabular}{@{}clccc@{}}
\toprule
$n$ & \textbf{Orbit Label} & $D_n$ & $\dim(\mathcal{O})$ & $h$\\
\midrule
1 & $A_1$ & 58 & 190 & 2\\
2 & $2A_1$ & 56 & 192 & 4\\
3 & $3A_1$ & 54 & 194 & 6\\
$\vdots$ & $\vdots$ & $\vdots$ & $\vdots$ & $\vdots$\\
12 & (EW anchor) & 36 & 212 & 24\\
15 & (Hadronic) & 30 & 218 & 30\\
$\vdots$ & $\vdots$ & $\vdots$ & $\vdots$ & $\vdots$\\
26 & $E_8(a_1)$ & 8 & 240 & 58\\
\bottomrule
\end{tabular}
\end{center}

\textbf{Full data:} Available in machine-readable format at the replication repository 
(see Section~\ref{app:repro}).

\subsection{Verification Algorithm}

The uniqueness is verified by:
\begin{enumerate}[leftmargin=1.5em, topsep=2pt, itemsep=1pt]
    \item Constructing the complete Hasse diagram of $\mathcal{N}(\mathfrak{e}_8)$ 
        from standard tables.
    \item Computing all cover relations with $\Delta D = 2$.
    \item Running dynamic programming to find \emph{all} paths from $D = 58$ to $D = 8$.
    \item Applying the cost function $\mathcal{S} = \sum_n (\Delta^2 \ln D_n)^2$ 
        (second-difference smoothness).
    \item Verifying that the minimum-cost path is unique (cost gap to second-best $= 0$ 
        indicates tie-breaker was not needed for the optimal path).
\end{enumerate}

%=============================================================================
\section{Boundary Trace Derivation}
\label{app:boundary}
%=============================================================================

This appendix derives the block index formula $k_B = \frac{3}{2} I_1(B)$ from 
the APS index theorem.

\subsection{Setup}

Consider a 6D manifold $M^6$ with boundary $\partial M^6 = \Sigma^5$. For a 
Weyl fermion coupled to a $U(1)$ gauge field $A$, the APS index theorem gives:
\begin{equation}
    \text{ind}(\not{\!D}) = \int_{M^6} \text{ch}(F) \wedge \hat{A}(M^6) 
    - \frac{1}{2}\eta(\not{\!D}_\Sigma)
\end{equation}
where $\eta$ is the $\eta$-invariant of the boundary Dirac operator.

\subsection{Flat Connection on Boundary Cycles}

On the orientable double cover $\widetilde{M}$, the boundary decomposes as:
\begin{equation}
    \partial\widetilde{M} = C_1 \cup C_2 \cup C_T
\end{equation}
(two boundary cycles plus the twist seam).

For a \emph{flat} $U(1)$ connection on each $C_i$ with holonomy $\theta_i = \oint_{C_i} A$, 
the $\eta$-invariant contribution becomes:
\begin{equation}
    \eta_i \propto \sum_{\text{fermions}} q^2 \cdot \theta_i
\end{equation}
where $q$ is the $U(1)$ charge.

\subsection{\texorpdfstring{Spectral Trace Coefficient and $\bone$}{Spectral Trace Coefficient and b1}}
The regulated determinant of the hypercharge-coupled Dirac operator contains a universal
UV logarithmic term whose coefficient is the quadratic charge trace:
\begin{equation}
    \ln\det(\not{D}_Y) \supset \frac{1}{2\pi}\left(\frac{5}{3}\sum_f Y_f^2\right)
    \ln\frac{\Lambda}{\mu}\;\int F^2.
\end{equation}
Thus $\bone = \frac{5}{3}\sum_f Y_f^2 = 41/10$ appears as a \emph{spectral index coefficient},
independent of IR thresholds, matching its role in the CFE.

\subsection{\texorpdfstring{Derivation of $k_B$}{Derivation of kB}}

\textbf{Step 1:} Each boundary cycle contributes a factor proportional to $I_1 = \sum q^2$.

\textbf{Step 2:} Three cycles contribute additively: factor 3.

\textbf{Step 3:} The double cover counts each physical state once across two sheets. 
The effective weight is halved: factor $1/2$.

\textbf{Result:}
\begin{equation}
    k_B = 3 \times \frac{1}{2} \times I_1(B) = \frac{3}{2} I_1(B)
\end{equation}

\textbf{Consistency check:} For the electroweak block, 
$I_1^{\text{EW}} = \frac{41}{48}$ (SM hypercharge sum), giving 
$k_{\text{EW}} = \frac{3}{2} \times \frac{41}{48} = \frac{41}{32}$, 
as stated in Table~\ref{tab:blocks}.

%=============================================================================
\section{\texorpdfstring{Seam Gluing and the $\eta$-Invariant Normalization}{Seam Gluing and the eta-Invariant Normalization}}
\label{app:seam}
%=============================================================================

This appendix provides the $\eta$-invariant gluing derivation used in Axiom~\ref{ax:GB}.

\subsection{Cut Space and Operators}
Let $\widetilde{M}$ be the orientable double cover of the M\"obius fiber and cut it along
the seam $\Gamma$ to obtain a manifold-with-boundary $\widetilde{M}_\Gamma$ with
boundary components $(C_1,C_2,\Gamma_+,\Gamma_-)$. Let $\not{D}_\Gamma$ be the Dirac operator
on $\widetilde{M}_\Gamma$ with APS boundary conditions on $C_{1,2}$ and a unitary matching
operator $U_\Gamma$ identifying $\Gamma_+$ with $\Gamma_-$.
On the seam we take the tangential 1D Dirac operator
\begin{equation}
    D_\Gamma = i\,\frac{d}{d\theta}, \qquad \theta\in[0,2\pi).
\end{equation}

\subsection{\texorpdfstring{Gluing Formula for the $\eta$-Invariant}{Gluing Formula for the eta-Invariant}}
The APS gluing theorem expresses the $\eta$-invariant of the glued manifold in terms of
the cut manifold plus a discrete spectral-flow correction:
\begin{equation}
    \eta(\not{D}_{\widetilde{M}}) = \eta(\not{D}_\Gamma) + 2\,\mathrm{SF}(U_\Gamma),
\end{equation}
where $\mathrm{SF}(U_\Gamma)\in\mathbb{Z}$ is the spectral flow of the matching operator.
In TFPT bookkeeping this contributes a quantized seam term
\begin{equation}
    \Delta_\Gamma = 2\pi\,\mathrm{SF}(U_\Gamma)
\end{equation}
to the curvature normalization.

\subsection{M\"obius Matching and the Minimal Gluing Class}
For the M\"obius $\mathbb{Z}_2$ identification, the matching satisfies $U_\Gamma^2=\mathbf{1}$.
The minimal nontrivial class is modeled by the unitary phase
\begin{equation}
    U_\Gamma(\theta)=e^{i\theta}.
\end{equation}

\subsection{Spectral Flow Computation}
Interpolate boundary conditions by a phase $\varphi$,
\begin{equation}
    \psi(\theta+2\pi)=e^{i\varphi}\psi(\theta), \qquad \varphi\in[0,2\pi],
\end{equation}
which gives the eigenvalues
\begin{equation}
    \lambda_n(\varphi)=n+\frac{\varphi}{2\pi}, \qquad n\in\mathbb{Z}.
\end{equation}
As $\varphi$ runs from $0$ to $2\pi$, exactly one eigenvalue crosses zero, so
$\mathrm{SF}(U_\Gamma)=1$ and $\Delta_\Gamma=2\pi$.

\begin{remark}
For unitary matching maps $U_\Gamma:S^1\to U(1)$, the spectral flow equals the winding
number of $\det U_\Gamma$, consistent with the $K^1(S^1)\cong\mathbb{Z}$ classification.
\end{remark}

\subsection{\texorpdfstring{Normalization of $\phitree$}{Normalization of phi_tree}}
The two physical boundary cycles contribute $2\pi+2\pi$, and the seam gluing adds
$\Delta_\Gamma=2\pi$, so the total curvature is $R_{\text{total}}=6\pi$. The stationarity
normalization $\varphi\cdot R_{\text{total}}=1$ therefore yields
\begin{equation}
    \phitree = \frac{1}{6\pi}.
\end{equation}

%=============================================================================
\section{\texorpdfstring{Derivation of the Coefficient 48 in $\deltatop$}{Derivation of the Coefficient 48 in delta_top}}
\label{app:48deriv}
%=============================================================================

This appendix explains the integer coefficient in $\deltatop = 48\cthree^4$.
The key point is that the canonical quantity is
\begin{equation}
    \deltatop = \frac{3}{256\pi^4},
\end{equation}
and the appearance of ``48'' is simply the same number expressed in units of $\cthree^4$.

\subsection{The Physical Setup}

The topological correction $\deltatop$ arises from the leading gauge-invariant correction 
to the tree-level geometric scale $\phitree$. The requirements are:
\begin{enumerate}[leftmargin=1.5em, topsep=2pt, itemsep=1pt]
    \item \textbf{Gauge invariance:} The correction must be built from $\cthree$ only 
        (no explicit $\alpha$ dependence at this order).
    \item \textbf{Dimensionlessness:} Since $\cthree$ is dimensionless, the correction 
        must be a pure power $\cthree^n$ with an integer coefficient.
    \item \textbf{Discreteness:} Topological invariants take discrete values; the 
        coefficient must be a rational number with simple structure.
\end{enumerate}

\subsection{\texorpdfstring{Power Counting: Why $\cthree^4$?}{Power Counting: Why c3^4?}}

The effective action on the M\"obius fiber includes the Chern--Simons term 
$\cthree\, a F\tilde{F}$. The leading topological correction comes from the 
one-loop determinant of the Dirac operator coupled to this background.

\textbf{Dimensional analysis:} The Chern--Simons coupling enters as $\cthree$. The 
Euler characteristic and Pontryagin invariants contribute at order $\cthree^2$ each. 
The product yields:
\begin{equation}
    \deltatop \sim \cthree^2 \times \cthree^2 = \cthree^4
\end{equation}

\textbf{No $\cthree^2$ term:} A correction of order $\cthree^2$ would require a 
pure Euler term, but this is already absorbed into $\phitree$ via Gauss--Bonnet. 
The next allowed correction is $\cthree^4$.

